% Nouveau EGA V (cote 25): archive cote 25, site page 126 (PDF page 127).
% English translation of the checked French annotated master.
% AI-translated by GPT-6 Astra (OpenAI, ChatGPT Pro mode).
\scanpage{25/126}
\mspage{\hand{29}}

\var{\margin{\unsure{by observing, by checking in §22 for the construction of the Grassmannian preliminaries}}}{\placedmargin{\unsure{by observing, by checking in §22 for the construction of the Grassmannian preliminaries}}}{marginal note placed}{ega25_25_126-v0001}
One can then express \var{\over{5}{4}}{4}{writing over}{ega25_25_126-v0002}.2. by saying that the set of classes
up to isomorphism of linear systems of divisors on X associated with
$\underline{L}$ , is in canonical bijective correspondence with the set \var{\xout{Grassx}}{}{typewriter strike-out}{ega25_25_126-v0003}
Grass(Q)(S\var{\xout{xxx}}{}{typewriter strike-out}{ega25_25_126-v0004} and, this correspondence being compatible with \unsure{the|of} arbitrary chang
of base, one sees that the functor
\[ S' \hand{\rightsquigarrow} \text{ensemble des classes (mod isomorphisme) de systèmes} \]
\[ \text{linéaires de diviseurs de } X_{S'}/S' \text{ associés à } \underline{L}_{S'} \]
is representable by the S-prescheme \uline{Grass} (Q).
\tnote{25-126-1}{The prose within the two displays is retained character for character with the French formulas. It reads: ``set of classes (mod isomorphism) of linear systems of divisors of'' the displayed base change, ``associated with'' the displayed invertible sheaf. Before ``changes of base'', the source has the candidates ``les'' and ``de''; both are represented by the marked grammatical alternatives. The source word ``changemen'' is clipped and is retained as the English fragment ``chang''.}

\var{\margin{Under the conditions of \unsure{4.1.|1.1.}, the lin. syst. of divisors is said to be \uline{complete} when \unsure{for} $\forall s\in S$, $E(s)\rightarrow H^0(X_s,\underline{L}|X_s)$ is \uline{bijective}. This is a condi-\unsure{tion} \uline{\unsure{A.R.}} \var{\deleted{\illegible{2}}}{}{deletion}{ega25_25_126-v0006} \unsure{(at least if} $g:X\rightarrow S$ \unsure{is of} finite pres.). if $g$ is proper, it implies that $g_*(\underline{L})$ is \uline{locally free of finite type}, and \unsure{u} is an \uline{isom.} \unsure{for} \illegible{3} a syst. exists.}}{\placedmargin{Under the conditions of \unsure{4.1.|1.1.}, the lin. syst. of divisors is said to be \uline{complete} when \unsure{for} $\forall s\in S$, $E(s)\rightarrow H^0(X_s,\underline{L}|X_s)$ is \uline{bijective}. This is a condi-\unsure{tion} \uline{\unsure{A.R.}} \var{\deleted{\illegible{2}}}{}{deletion}{ega25_25_126-v0006} \unsure{(at least if} $g:X\rightarrow S$ \unsure{is of} finite pres.). if $g$ is proper, it implies that $g_*(\underline{L})$ is \uline{locally free of finite type}, and \unsure{u} is an \uline{isom.} \unsure{for} \illegible{3} a syst. exists.}}{marginal note placed}{ega25_25_126-v0005}
% NOTE: Several left-margin line beginnings are clipped by the supplied images; unresolved fragments remain marked.
\var{\deleted{The following result, of a less technical nature than 5.2., should
rather precede the latter:

\uline{Corollary} 5.3. Let f:X \quad S be a faithfully flat morphism loc of finite pres
\var{\inserted[typed]{on X/S}}{on X/S}{insertion}{ega25_25_126-v0008}
. Then for every\var{\xout{xxfamily}}{}{typewriter strike-out}{ega25_25_126-v0009} linear system of divisors, parametrized
by a projective bundle P over S, the corresponding morphism
\[ \hand{\check{\ }}P \qquad \mathrm{Div}_{X/S} \]
is a monomorphism.

This means that two morphisms $S'$ \quad \hand{$\check{\ }$}P having the same composite with
the preceding morphism are equal. Making the base change $S'$ \quad S,
one is reduced \var{\xout{to showing that two}}{}{typewriter strike-out}{ega25_25_126-v0010} to the case of two sections of \hand{$\check{\ }$}P over S,
and}}{}{deletion}{ega25_25_126-v0007}

% NOTE: The paragraph from "It must also be made explicit" to "a) Suppose ... alge" is enclosed in a handwritten box.
It must also be made explicit in \var{\over{5}{4}}{4}{writing over}{ega25_25_126-v0011}.1. that $\underline{L}|X-Z$ is canonically
isomorphic to $f^{\hand{*}}(\underline{O}_P(1))$ (with the notation of the preceding no.), hence in
the case $D \hand{\in\emptyset}$ considered in \var{\over{\unsure{5|4}}{3}}{3}{writing over}{ega25_25_126-v0012}.3. , $\underline{L}|X-Z$ is precisely $\underline{L}_{X-Z}$
of \var{\over{\unsure{5|4}}{3}}{3}{writing over}{ega25_25_126-v0013}.3. b), and consequently $\underline{L}$ is precisely the canonical and
unique extension of $f^*(\underline{O}_P(1))$ to an invertible sheaf on X.

\uline{Proposition} \var{\over{5}{4}}{4}{writing over}{ega25_25_126-v0014}.3. Let D be a linear system of divisors on X/S
parametrized by \hand{$\check{\ }$}P , where g: X$\hand{\longrightarrow}$ S is \var{\deleted{\unsure{smooth}}}{}{deletion}{ega25_25_126-v0015} a flat morphism \var{\inserted{loc.}}{loc.}{insertion}{ega25_25_126-v0016} of finite presentation.

a) Suppose \var{\inserted{(g of finite pres. and}}{(g of finite pres. and}{insertion}{ega25_25_126-v0017} that for every $s \hand{\in} S$, denoting by k an alge

\var{\margin{Complete lin. \unsure{syst.} associated with $\underline{L}$ iff $g_*(\underline{L})$ is loc. free, its formation \unsure{commutes|commuting} with passage to the fibers, and for \unsure{every} $s\in S$, and $\varphi\in H^0(X_s,\underline{L}|X_s)$, $\varphi\neq0$, $\varphi$ is a regular section, condi}}{\placedmargin{Complete lin. \unsure{syst.} associated with $\underline{L}$ iff $g_*(\underline{L})$ is loc. free, its formation \unsure{commutes|commuting} with passage to the fibers, and for \unsure{every} $s\in S$, and $\varphi\in H^0(X_s,\underline{L}|X_s)$, $\varphi\neq0$, $\varphi$ is a regular section, condi}}{marginal note placed}{ega25_25_126-v0018}
% NOTE: The bottom handwritten note is cut off after condi at the right edge.
\tnote{25-126-2}{The abbreviations lin., syst., pres., loc. and isom. reproduce the corresponding abbreviations for linear, system, presentation, locally and isomorphism. The doubtful abbreviation A.R. remains unexpanded. The source's final ``algé'' continues onto the next page, and the bottom note ends with the clipped fragment ``condi''.}
